
\section{Method}

\begin{figure*}[t!]
  \centering
  \includegraphics[width=0.95\textwidth]{figs/mainfig_3.pdf}
  \caption{\textbf{\methodname pipeline.} Given editable superquadric controls with per-primitive spatial conditioning strengths $\tau_i$, SpaceFlow encodes the high-control regions and injects them into the TRELLIS sparse-structure flow process over a controlled time interval. The resulting structure is then passed to a texture flow model, where optional per-part text or image conditions are routed through PartField segments and refined with self-similarity guidance, producing a 3D asset with localized spatial and appearance control.
  \vspace{-4mm}}
  \label{fig:overview}
\end{figure*}

We present our method for part-wise spatial and appearance guidance for 3D Generation, detailing our superquadric primitives, pipeline architecture, and spatial and appearance control formulations.

\subsection{Problem Formulation}
% Given a text prompt $p$, a user-edited set of deformable superquadric primitives $\mathcal{S}=\{S_i\}_{i=1}^{N}$, per-primitive control levels $\boldsymbol{\tau}=\{\tau_i\}_{i=1}^{N}$, and optional appearance cues $\mathcal{A}=\{a_i\}$, our goal is to generate a textured 3D asset $X=(G,T)$ whose geometry and appearance obey the user constraints locally. Each primitive $S_i$ defines a spatial region of intent, while $\tau_i$ controls how strongly the generated geometry should follow that region. High values of $\tau_i$ enforce close adherence to the scaffold; lower values allow the generative prior to reinterpret or complete the region. When an appearance cue $a_i$ is provided, it should affect only the corresponding primitive region.
Given a text prompt $p$, a user-edited set of deformable superquadric primitives $\mathcal{S}=\{S_i\}_{i=1}^{N}$, per-primitive control levels $\boldsymbol{\tau}=\{\tau_i\}_{i=1}^{N}$ with $\tau_i \in \{\tau_{\text{low}}, \tau_{\text{high}}\}$, and optional appearance cues $\mathcal{A}=\{a_i\}$, our goal is to generate a textured 3D asset $X=(G,T)$ whose geometry and appearance obey the user controls locally. The conditioning strengths $\tau_{\text{low}}$ and $\tau_{\text{high}}$ are user-defined values within the range $[0, 12]$ satisfying $\tau_{\text{low}} < \tau_{\text{high}}$. Each primitive $S_i$ defines a spatial region of intent, while its control strength $\tau_i$ determines the local trade-off between shape fidelity and generative completion. Primitives assigned to the high-control strength $\tau_{\text{high}}$ enforce strict adherence to the input geometry, whereas primitives assigned to the low-control strength $\tau_{\text{low}}$ allow the generative prior to reinterpret or complete the region. When an appearance cue $a_i$ is provided, it must affect only the corresponding primitive region $i$.

% Unlike part-based generators, we do not require the final asset to be represented as a set of generated parts. The primitive scaffold is used only as localized guidance for a pretrained 3D generator: constrained regions should respect the user input, while unconstrained regions remain free to be completed plausibly.

\subsection{Superquadric Primitives}
\label{sec:superquadric_primitives}
Superquadrics are compact primitives that span ellipsoidal, cylindrical, rounded-box, and box-like shapes~\cite{1673799}. Each control primitive is parameterized as $S_i(\mathbf{R}_i,\mathbf{t}_i,\mathbf{s}_i,\boldsymbol{\epsilon}_i,\boldsymbol{\delta}_i)$, where $(\mathbf{R}_i,\mathbf{t}_i)$ is its pose (rotation matrix and translation vector), $\mathbf{s}_i = (s_{x,i}, s_{y,i}, s_{z,i})^\top$ its scale, $\boldsymbol{\epsilon}_i = (\epsilon_{1,i}, \epsilon_{2,i})^\top$ its shape exponents, and $\boldsymbol{\delta}_i$ its deformation parameters.
For a 3D point $\mathbf{x}$, let $\mathbf{u} = (u_x, u_y, u_z)^\top = \mathbf{R}_i^\top(\mathbf{x}-\mathbf{t}_i)$ be its coordinates in the primitive's local frame. Ignoring deformations, the superquadric volume is defined by the implicit inequality $F_i(\mathbf{x})\leq 1$, where the inside-outside function $F_i(\mathbf{x})$ is given by:
\vspace{-3mm}
\begin{equation}
F_i(\mathbf{x}) = \left( \left|\frac{u_x}{s_{x,i}}\right|^{2/\epsilon_{2,i}} + \left|\frac{u_y}{s_{y,i}}\right|^{2/\epsilon_{2,i}} \right)^{\epsilon_{2,i}/\epsilon_{1,i}} + \left|\frac{u_z}{s_{z,i}}\right|^{2/\epsilon_{1,i}} .
\label{eq:superquadric_implicit}
\end{equation}
In our scaffold, the deformation parameter $\boldsymbol{\delta}_i$ adds bending and tapering degrees of freedom, making the primitives more expressive than raw superquadrics. We use these deformable superquadrics as editable spatial control regions, not as the final output representation.


% \subsection{Two-Stage 3D Generation Backbone}
% \label{sec:trellis_backbone}

\subsection{Architecture Overview}

% Our method builds upon TRELLIS, which is a pretrained 3D generation backbone~\cite{Xiang_2025_CVPR}. TRELLIS represents an asset with a structured latent representation, SLAT, consisting of a sparse set of active voxels and local latent features attached to those voxels. The representation captures both geometry and appearance and can be decoded into meshes, radiance fields, or Gaussian splats.
Our method is built upon TRELLIS~\cite{Xiang_2025_CVPR}, a pretrained 3D generation backbone. TRELLIS represents 3D assets using a structured latent representation (SLAT), which consists of a sparse set of active voxels and local latent features attached to those voxels. This representation jointly captures geometry and appearance, and can be decoded into multiple formats including meshes, neural radiance fields, or Gaussian splats.

%TRELLIS generates SLAT in two stages. A first rectified-flow transformer generates the sparse voxel structure, which determines the coarse 3D support of the object. A second rectified-flow transformer then generates the local latent features on the active voxels, after which a decoder produces the final textured asset. 
TRELLIS generates this representation in two stages. First, a rectified-flow transformer generates the sparse voxel structure that defines the coarse 3D support of the object. Second, a subsequent rectified-flow transformer generates the local latent features on the active voxels, after which the VAE decoder produces the final textured asset. 

%This factorization makes TRELLIS well suited to our setting: we guide the structure stage with spatial conditions from user-edited primitives (see Section~\ref{sec:spatial_local_control}), and guide the latent stage with localized per-primitive appearance cues (see Section~\ref{sec:appearance_local_control}).

% An overview of our proposed pipeline is illustrated in Fig.~\ref{fig:overview}. The system takes as input a set of editable superquadric control shapes with associated spatial conditioning strengths $\tau_i$, along with optional text prompts or reference images for texture control. The spatial control inputs are processed via the Structure Flow Model (Section~\ref{sec:spatial_local_control}) to synthesize the coarse 3D voxel structure. The resulting structure is then decoded and passed to the Texture Flow Model (Section~\ref{sec:appearance_local_control}), where localized per-part text or image conditions are routed to specific sections of the asset via a segmented PartField representation and synthesized using self-similarity guidance.
This decoupled architecture is uniquely suited to our localized control setting. As outlined in Fig.~\ref{fig:overview}, the user-edited superquadrics and their spatial conditioning strengths $\tau_i$ guide the Structure Flow Model (Section~\ref{sec:spatial_local_control}) to generate the coarse 3D voxel structure. This structure is then decoded and passed to the Texture Flow Model (Section~\ref{sec:appearance_local_control}), where localized text prompts or reference images are routed to target sections of the asset via a segmented PartField representation and generated under self-similarity guidance.

\subsection{Local Spatial Control for Structure Generation}
\label{sec:spatial_local_control}
%describe the weighted avg,repaint, injection... @joan

To implement local spatial control during structure generation, we partition the user-defined superquadrics $\mathcal{S}$ into high-control ($\mathcal{S}_{\text{high}}$) and low-control ($\mathcal{S}_{\text{low}}$) subsets based on their individual conditioning strengths $\tau_i$. We also define noise thresholds $\tau_{\text{low}}$ and $\tau_{\text{high}}$ in the flow model's trajectory, which govern the timesteps over which these spatial controls are active. 

First, we render the complete set $\mathcal{S}$ and the high-control subset $\mathcal{S}_{\text{high}}$ independently. Both renders are encoded into the structure space using the VAE of TRELLIS to obtain the clean latents $\mathbf{z}_0^{\text{all}}$ and $\mathbf{z}_0^{\text{high}}$. Additionally, we define a binary spatial mask $\mathbf{M}_{\text{low}}$ by generating a bounding box around $\mathcal{S}_{\text{low}}$ with a $7\%$ spatial padding, which is then downsampled to the latent dimensions ($16 \times 16 \times 16$) using a $4 \times 4 \times 4$ 3D max pooling operation.

Following SpaceControl, we initialize the process by perturbing $\mathbf{z}_0^{\text{all}}$ with noise corresponding to the timestep $\tau_{\text{low}}$, yielding the initial noisy latent $\mathbf{z}_{\tau_{\text{low}}}^{\text{all}}$. We then begin the denoising process from this state to establish the coarse layout of the low-control components.

To ensure that the structural details of the high-control regions are preserved in the output, we inject the clean high-control latent $\mathbf{z}_0^{\text{high}}$ at each denoising step $t$. Using the low-control mask $\mathbf{M}_{\text{low}}$, we compute a spatially masked weighted average to obtain the updated latent $\mathbf{z}'_t$:
\begin{equation}
\mathbf{z}'_t = (1 - \mathbf{M}_{\text{low}}) \odot \left( \alpha \mathbf{z}_0^{\text{high}} + (1 - \alpha) \mathbf{z}_t \right) + \mathbf{M}_{\text{low}} \odot \mathbf{z}_t
\label{eq:high_control_injection}
\end{equation}
where $\mathbf{z}_t$ is the model's denoised prediction, $\odot$ denotes the element-wise product, and the blending weight is set to $\alpha = 0.18$. This hyperparameter was chosen experimentally but can be adjusted depending on the specific asset being generated.

To improve the transition boundaries between the high, and low, control regions, we apply a resampling strategy following RePaint~\cite{lugmayr2022repaint}. Following the noise scheduler, we perform $N$ iterations of a resampling loop. In each iteration, we first perturb the blended sample from $t_{i-1}$ back to step $t_i$. We then denoise the sample back to $t_{i-1}$ using a single solver step and re-apply the blending formulation in Eq.~\eqref{eq:high_control_injection}. This iterative blending is performed for the steps between $\tau_{\text{low}}$ and $\tau_{\text{high}}$, corresponding to the \textit{Conditioned Flow} phase shown in Fig.~\ref{fig:overview}.

% To improve the transition boundaries between the high- and low-control regions, we apply a resampling strategy inspired by RePaint~\cite{repaint}. We define the timestep difference $\Delta t = t_i - t_{i-1}$ following the noise scheduler and perform $N$ iterations of a resampling loop. In each iteration, we first perturb the blended sample back to step $t_i$. 

% Following the linear interpolation framework of flow matching, this local re-noising step over the interval $\Delta t$ is formulated as:
% \begin{equation}
% \mathbf{z}_{t_i} = (1 - \Delta t) \mathbf{z}'_{t_{i-1}} + \Delta t \mathbf{\epsilon}, \quad \mathbf{\epsilon} \sim \mathcal{N}(\mathbf{0}, \mathbf{I})
% \label{eq:repaint_renoising}
% \end{equation}
%where $\mathbf{z}'_{t_{i-1}}$ acts as the local clean reference. 



\subsection{Matching Generated Parts to Superquadrics}
\label{sec:part_matching}
To route appearance cues to the correct regions, we establish a correspondence between the generated structure and the input primitives $\mathcal{S}$. A direct geometric assignment is unreliable: low-control regions are reshaped by the generative prior and no longer align with their primitives, so nearest-primitive labelling yields fragmented, semantically inconsistent regions. We therefore match parts in a learned part-aware feature space.

Given the generated structure mesh, we extract a PartField~\cite{liu2025partfield} triplane feature field $\mathcal{P}$, which encodes part-level semantics learned across large 3D collections. For each active structure voxel $\mathbf{v}_j$ on the $64^3$ grid, we sample $\mathcal{P}$ at its normalized position by interpolating the three axis-aligned planes and summing them, yielding a $448$-dimensional descriptor $\mathbf{f}_j$. We $\ell_2$-normalize these descriptors and cluster them with $K$-means into $K{=}10$ semantic parts, assigning each voxel a part label $\ell_j$.

We then assign each part-cluster to a single primitive. Each $S_i$ is voxelized to its set of surface cells $\mathcal{B}_i$. For each cluster $C_k$ we take the voxel $\mathbf{c}_k$ closest to the cluster's spatial mean as its representative, and assign the whole cluster to the nearest primitive:
\vspace{-2mm}
\begin{equation}
a(k) = \arg\min_{i} \min_{\mathbf{b}\in\mathcal{B}_i} \lVert \mathbf{c}_k - \mathbf{b}\rVert_2 .
\label{eq:cluster_assignment}
\end{equation}
Every voxel inherits the primitive label $a(\ell_j)$ of its cluster. Operating at the cluster level rather than per voxel makes the assignment robust: an entire semantic part is routed to one primitive even where the generated geometry has drifted from the scaffold. The resulting labels are downsampled to the $32^3$ latent grid to form a dense routing volume $\mathbf{R}$, whose occupied cells store the index of the assigned primitive ($0$ denotes the global fallback).

Because routing is performed on the coarse $32^3$ grid, thin or small primitives (\eg wheels) may collapse to too few cells to carry their local cue. We therefore (i) fall back to a purely geometric per-voxel nearest-primitive assignment whenever cluster-based routing fails to cover a primitive holding an active local cue, and (ii) adaptively grow any under-covered primitive up to a minimum cell count by dilation, never overwriting cells already claimed by another primitive.

\subsection{Part-Wise Appearance Guidance}
\label{sec:appearance_local_control}
Given the matched parts, the second flow stage of TRELLIS synthesizes the SLAT latent features that determine the asset's appearance. We adopt the optimization-guided rectified-flow scheme of GuideFlow3D~\cite{sayandsarkar_2025_guideflow3d}, which interleaves flow sampling with a self-similarity guidance loss, and extend it with part-wise conditioning.

\noindent\textbf{Local condition routing.}
Each appearance cue (text or image) is encoded by TRELLIS into a conditioning embedding. We collect a global embedding $\mathbf{e}_0$ from the global prompt, together with one embedding $\mathbf{e}_i$ per primitive carrying a local cue. Using the routing volume $\mathbf{R}$ (Sec.~\ref{sec:part_matching}), each latent voxel is mapped to one such embedding. Inside every cross-attention layer of the flow model, a voxel attends only to its routed embedding: queries are evaluated against all candidate embeddings and, for each voxel, the output corresponding to its assigned condition is selected, with unrouted voxels defaulting to $\mathbf{e}_0$. This confines each appearance specification to the voxels of its target part and prevents texture leakage into the rest of the asset.

\noindent\textbf{Self-similarity guidance.}
To keep the synthesized appearance coherent within parts, we guide the latents with a supervised contrastive self-similarity loss that uses the PartField clusters $\{\ell_j\}$ as supervision. Voxels sharing a part label are attracted in feature space while those in different parts are repelled:
\vspace{-1mm}
\begin{equation}
\mathcal{L} = -\frac{1}{|\mathcal{V}|}\sum_{j} \log \frac{\sum_{k:\,\ell_k=\ell_j,\,k\neq j}\exp(s_{jk})}{\sum_{k\neq j}\exp(s_{jk})},
\label{eq:contrastive}
\end{equation}
where $\mathcal{V}$ is the set of structure voxels and $s_{jk}$ is the cosine similarity between the latent features of voxels $j$ and $k$. The loss is evaluated in chunks to avoid materializing the full $|\mathcal{V}|\times|\mathcal{V}|$ similarity matrix.

\noindent\textbf{Optimization.}
We optimize the per-voxel latent features over the rectified-flow trajectory. At each step we (i) take a flow sampler step under the routed conditions to move the latents toward the data manifold, and (ii) update them with one AdamW step on $\mathcal{L}$, keeping the lowest-loss iterate. The optimized latents are finally denormalized and decoded by TRELLIS into the textured 3D asset. Setting the guidance weight to zero recovers pure part-routed conditional sampling, isolating the effect of the self-similarity term.


%\subsection{Implementation Details}



% \begin{figure}
%     \centering
%     \includegraphics[width=1\linewidth]{figs/condit-flow-iter.pdf}
%     \caption{}
%     \label{fig:condit-flow-iter}
% \end{figure}

% \begin{figure}
%     \centering
%     \includegraphics[width=1\linewidth]{figs/condit-routing-outline.pdf}
%     \caption{}
%     \label{fig:condit-routing-outline}
% \end{figure}

% \begin{figure}
%     \centering
%     \includegraphics[width=1\linewidth]{figs/condit-routing-details.pdf}
%     \caption{}
%     \label{fig:condit-routing details}
% \end{figure}

